\section{Tax information before tax calculation}
\label{sec:shared-tax-information}

The mixed book has recorded profit, but it still lacks a tax assessment.
The earlier return-only reference does contain a tax-close method. Whether
that method can be reused is a substantive question: the new book contains
lease rights and a modification gain that the old book did not represent.
The unchanged method refuses the mixed state because its closed state shape
does not admit those additional accounts and lease records. The two original
tax examples still reproduce their complete results. Thus the refusal identifies
an incompatible extension, rather than a regression in the original method.

\subsection{A new income category changes the tax starting point}

The original pretax function sums the original income accounts, excluding
current and deferred tax. If that function is applied diagnostically to the
mixed balances, it returns a loss of 15. The complete recorded report instead
includes the lease-modification gain of 20. Under the signed-account convention,
the additional \texttt{other\_gains} balance is $-20$, so
\[
  P_{\mathrm{old\ accounts}}=-15,
  \qquad
  P_{\mathrm{mixed}}=-15-(-20)=5.
\]
The difference is an omitted income category, not a rounding error. Replacing
the state validator while leaving the original income list unchanged would
therefore calculate the wrong pretax starting point for the stated mixed-book
target. It would not preserve the meaning of the old calculation.

Even the corrected starting point does not determine tax. The return example
has a signed reconciliation adjustment of $-24$. Mechanical reuse of that
adjustment and the original illustrative rate would give
\[
  \frac{-15-24}{4}=-\frac{39}{4},
  \qquad
  \frac{5-24}{4}=-\frac{19}{4}.
\]
These expressions diagnose the effect of the omitted gain; neither is an
assessed tax or a cash refund. The lease deduction timing, the tax treatment
of the modification, the remaining right-of-use tax bases and the treatment
of a loss have not been supplied. The original finite method refuses its
negative-current-tax route. Copying the rate or reversing every deferred-tax
balance would not supply the missing facts.

Dates create a separate incompatibility. The original close on 31 December
stores 1 January as an exclusive tax endpoint. The mixed-book invariant
currently requires that endpoint to be no later than the state date. Those
two meanings cannot be combined by changing only the endpoint field. Moreover,
the mixed example's execution window runs from 1 January 2024 to 1 January 2026:
731 days. That window contains the selected lease and return events, but does
not establish a single annual tax interval or complete fiscal activity.

\subsection{Keep missing declarations and rejected premises visible}

A separate information assessor now accepts the unchanged recorded-report
request, an explicitly stated tax interval and identified information
declarations. The interval distinguishes its start, its last included calendar
day and its exclusive endpoint. The first finite implementation requires the
start to equal the recorded opening and every requested event to lie inside
the interval. The last recorded event may precede the asserted coverage date;
that assertion remains a declaration, not evidence that all activity is known.

For this mixed example, the assessor identifies 61 required declarations. They
cover jurisdiction and regime, the interval basis and coverage findings, rate,
loss and cash-settlement treatment, deferred recognition and recovery, every
pretax account's treatment and signed reconciliation adjustment, balance-sheet
tax bases, and the tax bases and timing of each lease right. Zero balances
remain in the account population. Both lease rights A and B remain relevant:
the native state retains fully terminated B with zero asset and liability,
and its earlier modification still requires a stated treatment.

Each supplied value carries an assumption or source-reference description.
Each unknown carries its reason. A source reference is retained as supplied;
the assessor does not authenticate it. With no information supplied, all
61 declarations remain unresolved. Removing any one from the hypothetical
complete control again leaves an unresolved requirement. A false coverage or
recognition finding remains false and is reported as unsupported. Extra fields,
wrong scalar types, forged request identities and incompatible interval
boundaries are invalid.

The complete native report accompanies the assessment. Its stopped event
prefix, component memory, gross cash movements and unresolved annual coverage
are preserved. When a requested payment fails, supplying every information
declaration does not erase that failure: completion of the requested sequence
remains unresolved. A first-year interval can be stated separately when its
requested events fit that interval, without relabelling the two-year example.

\subsection{Supplied information is not an executable tax rule}

The positive result is deliberately called \texttt{information\_supplied}.
It means the finite declaration set is present and has the required types.
The hypothetical controls exercise that representation; their supplied zero
bases and adjustments are not inferred from the financial statements. A text
field describing a deduction can still fail to specify a valid calculation.
Likewise, individually typed amounts can be inconsistent when aggregated.

Accordingly, every result keeps calculation readiness, source and accounting
admission, tax execution and annual-close readiness false. No tax amount is
posted. The next calculation design must state how account adjustments,
right-level tax bases, deferred recognition and settlements relate, and must
check their joint consistency. This preserves useful progress: missing facts
and competing declarations can be represented and examined now, while the
remaining numerical and accounting questions remain explicit.
